% Member 2 report sections: order fulfilment (HLD 5 services 4, 5 and 11; HLD 6.1; one HLD 13 trade-off).
% Body only, so the team report can % Member 2 report sections: order fulfilment (HLD 5 services 4, 5 and 11; HLD 6.1; one HLD 13 trade-off).
% Body only, so the team report can % Member 2 report sections: order fulfilment (HLD 5 services 4, 5 and 11; HLD 6.1; one HLD 13 trade-off).
% Body only, so the team report can % Member 2 report sections: order fulfilment (HLD 5 services 4, 5 and 11; HLD 6.1; one HLD 13 trade-off).
% Body only, so the team report can \input{member2-fulfilment}. Needs: array, booktabs, tikz (arrows.meta,
% positioning, calc).
% member2-standalone.tex compiles these sections on their own.

\section{Fulfilment Orchestrator (HLD 5, service 4)}

The Fulfilment Orchestrator turns each paid order into a job and keeps it moving until a rider takes the bag, so that
no order gets stuck or is acted on twice. It listens for \texttt{order.paid}, creates a \texttt{FulfilmentJob} with
one pick task per bin (oldest batch first, in walking order) and owns the job's state.

Every state change is one conditional \texttt{UPDATE} that names the state it expects:
\begin{center}
  \texttt{UPDATE fulfilment\_jobs SET state = :new WHERE id = :job AND state = :expected}
\end{center}
Of two requests racing for the same change, exactly one wins and the other is told what happened. This covers
claiming, cancelling, scanning, packing, stall release and handover.

\begin{figure}[ht]
  \centering
  \begin{tikzpicture}[
      >={Stealth[length=2.2mm]}, node distance=17mm,
      state/.style={draw=black!45, rounded corners=3pt, minimum width=24mm, minimum height=10mm, font=\bfseries},
      sale/.style={state, draw=blue!70!black, fill=blue!8, line width=0.9pt},
      edgelabel/.style={font=\small, text=black!65}]
    \node[state] (created) {created};
    \node[state, right=of created] (picking) {picking};
    \node[sale, right=of picking] (packed) {packed};
    \node[state, right=of packed] (handed) {handed over};
    \node[state, below=16mm of $(created)!0.5!(picking)$] (cancelled) {cancelled};
    \node[edgelabel, below=1mm of packed] {sale committed here};

    \draw[->, black!55] (created) -- node[edgelabel, above] {claim} (picking);
    \draw[->, black!55] (picking) -- node[edgelabel, above] {pack} (packed);
    \draw[->, black!55] (packed) -- node[edgelabel, above] {scan bag} (handed);
    \draw[->, black!55] (picking.north) -- ++(0, 7mm) -| node[edgelabel, pos=0.25, above]
      {idle 5 min, nothing in tote} (created.north);
    \draw[->, black!55] (created.south) |- node[edgelabel, pos=0.3, left] {cancel} (cancelled.west);
    \draw[->, black!55] (picking.south) |- node[edgelabel, pos=0.3, right] {cancel, or nothing to pack}
      (cancelled.east);
  \end{tikzpicture}
  \caption{Fulfilment job lifecycle. Stock leaves the books only at packing; every earlier exit only releases holds.}
  \label{fig:job-lifecycle}
\end{figure}

Figure~\ref{fig:job-lifecycle} shows the lifecycle. Before packing, an order can leave the line by a cancel, or go back
to the queue when its picker goes idle; both only release holds. After packing, the only way on is to its own rider.

\begin{itemize}
  \item \textbf{Store load.} More than 5 orders waiting marks the store \texttt{busy}, and serviceability adds
    5 minutes to the ETA.
  \item \textbf{Cancel after payment.} Allowed while the order is \texttt{created} or \texttt{picking}: its holds are
    given back and nothing leaves on-hand stock. A packed order would be a return, which is out of scope.
  \item \textbf{Stuck orders.} No claim, scan or shelf report for 5 minutes: with nothing in the tote, the order goes
    back to the queue (checked every 60\,s); with items in a tote, it is flagged for a person. Packed bags waiting
    over 5 minutes for a rider are flagged too.
  \item \textbf{KPIs} over the last 30 days: order-to-packed p50 and p90, claim-to-packed p90, packed-to-handover
    p90, fill rate (units picked $\div$ units ordered on the reservation) and items per picker-hour (a picker's
    overlapping orders counted once).
\end{itemize}

Percentiles take the value at index $\lfloor q \times (n - 1) \rfloor$ of the sorted list, the nearest rank below.
With few orders this errs low: the p90 of 30, 60 and 90\,s is 60\,s.

\section{Pick \& Pack (HLD 5, service 5)}

Pick \& Pack gives pickers work, checks every scan and packs the order. Packing is the moment the sale leaves the
books; until then the units are only held, so a cancel never has to put stock back.

\begin{itemize}
  \item \textbf{Claiming.} Off-peak, a picker gets the oldest waiting order. With 4 or more waiting, they get the
    3 oldest, each in its own tote (A, B, C), with one route through all their lines in walking order.
  \item \textbf{Scanning.} Every unit is barcode-checked as it goes into the tote (``3 of 6 scanned''). A wrong item,
    or one unit more than the line needs, changes nothing.
  \item \textbf{Empty shelf.} The picker says how many were there, and those stay picked. For the rest, the count is
    corrected (a \texttt{not\_found} ledger row) and the line moves to another bin with spare stock, never units
    another customer holds. With none left it is dropped: the hold is released, \texttt{order.line\_dropped} carries
    the refund, and a substitute is offered (same category, enough to sell, no dearer than the missing item, nearest
    price). The rest of the bin is written off too, so customers cannot buy stock that is not there. Units another
    unpacked order is waiting for stay, so held units never exceed those on hand, until that order's picker reports
    the same empty bin.
  \item \textbf{Packing.} Commits the sale for every picked line (\texttt{qty\_on\_hand} and \texttt{qty\_reserved}
    both drop) and publishes \texttt{order.ready\_for\_pickup}. An order with nothing picked is cancelled instead.
  \item \textbf{The handheld.} \texttt{/picker.html} is a phone screen: one line at a time (``Go to bin A-01-04,
    pick 2 $\times$ Basmati rice''), big buttons, and the tote shown at peak. A handheld that sends its picker id
    cannot act on an order that has since gone to someone else.
\end{itemize}

\section{Rider Handoff (HLD 5, service 11)}

The Rider Handoff Adapter makes sure the right bag goes to the right rider. Before it, handover was one unchecked
call, and a rider could leave with someone else's order.

\begin{enumerate}
  \item \textbf{Assigned.} When an order is packed, the rider system sends a rider. In the prototype a stand-in picks
    one from a fixed list (R-11 to R-15) inside the packing transaction, so a packed order always has a rider.
  \item \textbf{Arrived.} The rider checks in at the desk for an order and is told its crate slot. Another rider's
    order (\texttt{wrong\_rider}), or one not yet packed (\texttt{not\_ready}), is refused.
  \item \textbf{Collected.} The desk scans the bag before handing it over. A bag from another order
    (\texttt{wrong\_bag}) or a rider other than the assigned one (\texttt{wrong\_rider}) is refused and the order
    stays packed; the refusal names the crate slot holding the right bag.
\end{enumerate}

The time from packed to handed over is a fulfilment KPI (p90), so slow pick-ups show up next to slow packing.

\section{Order to dispatch (HLD 6.1)}

A paid order reaches a rider in seven steps, and stock leaves the books only at packing. That one rule keeps every
failure path cheap: before packing, any way out only gives a hold back.

\begin{enumerate}
  \item \textbf{Hold.} At checkout, Inventory holds the whole cart for 10 minutes, or none of it.
  \item \textbf{Pay.} Payment confirms the hold and publishes \texttt{order.paid}. The orchestrator creates the job
    and its pick tasks, oldest batch first, in walking order.
  \item \textbf{Claim.} A picker claims one order, or three at peak, with a conditional \texttt{UPDATE}.
  \item \textbf{Pick.} The picker walks the route and scans each unit into its tote.
  \item \textbf{Pack.} At the desk the sale is committed and a rider is assigned.
  \item \textbf{Check in.} The rider is told which crate slot holds the bag.
  \item \textbf{Hand over.} The desk scans the bag; only this order's bag goes to this order's rider.
\end{enumerate}

\paragraph{When an item is missing} (the HLD 9.2 saga: next bin, else substitute, else refund). Units the picker
found stay picked. The rest move to another bin with spare stock if there is one. Otherwise they are dropped and
refunded through \texttt{order.line\_dropped}, and a substitute is offered.

To accept a substitute, the demo has the picker approve it on the customer's behalf; a real app would ask the
customer on their phone. Only the item on offer can be accepted, and it is never dearer than the missing one, so
whoever accepts can never make the customer pay more than they were refunded. Only units the order really dropped
can be replaced, and they are marked so they cannot be
replaced twice. The substitute is held in the same transaction that changes the order line
(\texttt{inventory.swap\_line}, from Member~1); if it sold out since the offer, nothing changes. It joins the walk as
a new line, and \texttt{order.line\_substituted} charges for it, since the dropped units were already refunded.

\paragraph{Cancel and refunds.} Until packing, a cancel gives back every hold, and \texttt{order.cancelled} refunds
only those units. Units dropped earlier were refunded by their own \texttt{order.line\_dropped}, so the two never add up
to more than the order: in the tests, Rs~400 dropped plus Rs~560 cancelled is the order's Rs~960. After packing, a
cancel would be a return, which is out of scope, so it is refused (\texttt{too\_late}).

\section{Measurements}

Batch picking halves the walking at peak. The timing run's times are compressed, so they compare batching with
one order per walk; they do not test the HLD 12 time targets.

\paragraph{Walking distance} (\texttt{scripts/compare\_picking.py}, task 2.6). The same 30 random paid orders were
picked twice. One order per walk took 30 walks and 622 steps; batches of up to 3 took 12 walks and 322 steps, 48\%
less (46--49\% on other random seeds). A walk's distance is twice the furthest bin it visits, counted in bin steps.

\paragraph{Timing run} (\texttt{scripts/simulate\_orders.py}, task 2.8). 100 random paid orders, all waiting at the
start, were picked and packed through the real API, with one walk step taking 0.05\,s (seed 7). The stock audit was
clean after every run (Table~\ref{tab:timing-run}).

\begin{table}[ht]
  \centering
  \small
  \begin{tabular}{@{}lrrrl@{}}
    \toprule
    KPI & \shortstack[r]{One order per walk\\1 picker} & \shortstack[r]{Batches\\1 picker}
        & \shortstack[r]{Batches\\3 pickers} & HLD 12 target \\
    \midrule
    Walks                       & 100      & 34       & 34       &               \\
    Time to clear the queue (s) & 113      & 55       & 19       &               \\
    Order to packed, p90 (s)    & 103.3    & 49.6     & 17.0     & $\le 120$     \\
    Order to packed, p50 (s)    & 58.1     & 28.6     & 9.9      &               \\
    Claim to packed, p90 (s)    & 1.8      & 1.9      & 1.9      &               \\
    Fill rate                   & 100\%    & 100\%    & 100\%    & $\ge 99\%$    \\
    Items per picker-hour       & 8{,}905  & 18{,}149 & 18{,}278 & $\ge 100$     \\
    \bottomrule
  \end{tabular}
  \caption{Timing run: 100 paid orders, 0.05\,s per walk step, seed 7. The time-based targets are shown for
    reference only: with walking compressed, those rows are not a pass.}
  \label{tab:timing-run}
\end{table}

\paragraph{What the simulation can and cannot tell us.} Its comparisons hold, because every run shares the same orders,
store layout and walking speed: batching cut walks from 100 to 34 and halved the time to clear the queue. It also
shows the code stays correct under load, since every order was packed and the stock audit stayed clean.

Its absolute numbers do not hold. Walking is compressed, and scanning, packing and handover take no time, so 18,000
items per picker-hour describes the code, not people. All orders arrive at once, so order-to-packed mostly measures how
fast the pickers drain the queue.

Fill rate is 100\% only because no shelf is ever empty in the simulation. Real numbers need timed pickers in a real or
mock-up store, with orders arriving through the day.

\section{Trade-off: batch picking at peak (HLD 13)}

We batch orders only when at least 4 are waiting. Batching halves the walking at peak, and the threshold keeps it from
making an order wait when the store is quiet.

HLD 13.1 chose ``single-order off-peak, batch picking at peak'' because ``batching raises throughput at peak but adds
wait when orders are few''. We measured the first half: 48\% less walking for 30 orders. With one picker, a 100-order
queue cleared in 55\,s instead of 113\,s (compressed time).

The second half is handled by design rather than measured. With fewer than 4 orders waiting, a picker takes one at
once, so no order waits for a batch to fill. With 4 or more, a batch of 3 still leaves at least one order for the next
picker.

What batching still costs:
\begin{itemize}
  \item Each order waits while the picker walks to the other orders' bins. In the timing run, claim-to-packed p90
    barely moved (1.8\,s to 1.9\,s), but on a real floor this is real walking time.
  \item The picker must put each item in the right tote, and nothing checks the tote yet.
  \item A batch is the 3 oldest orders, not 3 whose items sit close together, so some batches save little walking.
\end{itemize}

\section{Limits and open items}

The fulfilment prototype's limits are shortcuts, each with a clear next step (Table~\ref{tab:limits}).

\begin{table}[!htb]
  \centering
  \footnotesize
  \begin{tabular}{@{}>{\raggedright\arraybackslash}p{0.27\linewidth}>{\raggedright\arraybackslash}p{0.36\linewidth}
                  >{\raggedright\arraybackslash}p{0.31\linewidth}@{}}
    \toprule
    Limit & What it means & Next step \\
    \midrule
    The picker accepts a substitute for the customer
      & The customer is not asked; the picker decides on their behalf, but only for the offered item, never a dearer one
      & Ask the customer on their phone, with a timeout \\
    \addlinespace[3pt]
    Units already in a tote when their batch expires
      & They are still packed and delivered; only the rest of the batch is refused at the scanner
      & Check the expiry date again at packing \\
    \addlinespace[3pt]
    Units typed as ``found'' are not scanned
      & They skip the barcode check
      & Scan them before packing \\
    \addlinespace[3pt]
    The picker check is opt-in
      & A handheld that sends no \texttt{picker\_id} can still act on an order since given to someone else
      & The demo station and phone screen send it; make it required \\
    \addlinespace[3pt]
    Riders are made up
      & A fixed list stands in for the rider system, and the desk checks riders in
      & Connect a real rider system \\
    \addlinespace[3pt]
    Refunds use today's price
      & A price change between checkout and refund would refund the wrong amount
      & Store the price on each reservation line \\
    \addlinespace[3pt]
    Batches are the oldest orders
      & Orders are not grouped by where their items sit
      & Group by zone at peak \\
    \bottomrule
  \end{tabular}
  \caption{Limits of the fulfilment prototype and what would close them.}
  \label{tab:limits}
\end{table}
. Needs: array, booktabs, tikz (arrows.meta,
% positioning, calc).
% member2-standalone.tex compiles these sections on their own.

\section{Fulfilment Orchestrator (HLD 5, service 4)}

The Fulfilment Orchestrator turns each paid order into a job and keeps it moving until a rider takes the bag, so that
no order gets stuck or is acted on twice. It listens for \texttt{order.paid}, creates a \texttt{FulfilmentJob} with
one pick task per bin (oldest batch first, in walking order) and owns the job's state.

Every state change is one conditional \texttt{UPDATE} that names the state it expects:
\begin{center}
  \texttt{UPDATE fulfilment\_jobs SET state = :new WHERE id = :job AND state = :expected}
\end{center}
Of two requests racing for the same change, exactly one wins and the other is told what happened. This covers
claiming, cancelling, scanning, packing, stall release and handover.

\begin{figure}[ht]
  \centering
  \begin{tikzpicture}[
      >={Stealth[length=2.2mm]}, node distance=17mm,
      state/.style={draw=black!45, rounded corners=3pt, minimum width=24mm, minimum height=10mm, font=\bfseries},
      sale/.style={state, draw=blue!70!black, fill=blue!8, line width=0.9pt},
      edgelabel/.style={font=\small, text=black!65}]
    \node[state] (created) {created};
    \node[state, right=of created] (picking) {picking};
    \node[sale, right=of picking] (packed) {packed};
    \node[state, right=of packed] (handed) {handed over};
    \node[state, below=16mm of $(created)!0.5!(picking)$] (cancelled) {cancelled};
    \node[edgelabel, below=1mm of packed] {sale committed here};

    \draw[->, black!55] (created) -- node[edgelabel, above] {claim} (picking);
    \draw[->, black!55] (picking) -- node[edgelabel, above] {pack} (packed);
    \draw[->, black!55] (packed) -- node[edgelabel, above] {scan bag} (handed);
    \draw[->, black!55] (picking.north) -- ++(0, 7mm) -| node[edgelabel, pos=0.25, above]
      {idle 5 min, nothing in tote} (created.north);
    \draw[->, black!55] (created.south) |- node[edgelabel, pos=0.3, left] {cancel} (cancelled.west);
    \draw[->, black!55] (picking.south) |- node[edgelabel, pos=0.3, right] {cancel, or nothing to pack}
      (cancelled.east);
  \end{tikzpicture}
  \caption{Fulfilment job lifecycle. Stock leaves the books only at packing; every earlier exit only releases holds.}
  \label{fig:job-lifecycle}
\end{figure}

Figure~\ref{fig:job-lifecycle} shows the lifecycle. Before packing, an order can leave the line by a cancel, or go back
to the queue when its picker goes idle; both only release holds. After packing, the only way on is to its own rider.

\begin{itemize}
  \item \textbf{Store load.} More than 5 orders waiting marks the store \texttt{busy}, and serviceability adds
    5 minutes to the ETA.
  \item \textbf{Cancel after payment.} Allowed while the order is \texttt{created} or \texttt{picking}: its holds are
    given back and nothing leaves on-hand stock. A packed order would be a return, which is out of scope.
  \item \textbf{Stuck orders.} No claim, scan or shelf report for 5 minutes: with nothing in the tote, the order goes
    back to the queue (checked every 60\,s); with items in a tote, it is flagged for a person. Packed bags waiting
    over 5 minutes for a rider are flagged too.
  \item \textbf{KPIs} over the last 30 days: order-to-packed p50 and p90, claim-to-packed p90, packed-to-handover
    p90, fill rate (units picked $\div$ units ordered on the reservation) and items per picker-hour (a picker's
    overlapping orders counted once).
\end{itemize}

Percentiles take the value at index $\lfloor q \times (n - 1) \rfloor$ of the sorted list, the nearest rank below.
With few orders this errs low: the p90 of 30, 60 and 90\,s is 60\,s.

\section{Pick \& Pack (HLD 5, service 5)}

Pick \& Pack gives pickers work, checks every scan and packs the order. Packing is the moment the sale leaves the
books; until then the units are only held, so a cancel never has to put stock back.

\begin{itemize}
  \item \textbf{Claiming.} Off-peak, a picker gets the oldest waiting order. With 4 or more waiting, they get the
    3 oldest, each in its own tote (A, B, C), with one route through all their lines in walking order.
  \item \textbf{Scanning.} Every unit is barcode-checked as it goes into the tote (``3 of 6 scanned''). A wrong item,
    or one unit more than the line needs, changes nothing.
  \item \textbf{Empty shelf.} The picker says how many were there, and those stay picked. For the rest, the count is
    corrected (a \texttt{not\_found} ledger row) and the line moves to another bin with spare stock, never units
    another customer holds. With none left it is dropped: the hold is released, \texttt{order.line\_dropped} carries
    the refund, and a substitute is offered (same category, enough to sell, no dearer than the missing item, nearest
    price). The rest of the bin is written off too, so customers cannot buy stock that is not there. Units another
    unpacked order is waiting for stay, so held units never exceed those on hand, until that order's picker reports
    the same empty bin.
  \item \textbf{Packing.} Commits the sale for every picked line (\texttt{qty\_on\_hand} and \texttt{qty\_reserved}
    both drop) and publishes \texttt{order.ready\_for\_pickup}. An order with nothing picked is cancelled instead.
  \item \textbf{The handheld.} \texttt{/picker.html} is a phone screen: one line at a time (``Go to bin A-01-04,
    pick 2 $\times$ Basmati rice''), big buttons, and the tote shown at peak. A handheld that sends its picker id
    cannot act on an order that has since gone to someone else.
\end{itemize}

\section{Rider Handoff (HLD 5, service 11)}

The Rider Handoff Adapter makes sure the right bag goes to the right rider. Before it, handover was one unchecked
call, and a rider could leave with someone else's order.

\begin{enumerate}
  \item \textbf{Assigned.} When an order is packed, the rider system sends a rider. In the prototype a stand-in picks
    one from a fixed list (R-11 to R-15) inside the packing transaction, so a packed order always has a rider.
  \item \textbf{Arrived.} The rider checks in at the desk for an order and is told its crate slot. Another rider's
    order (\texttt{wrong\_rider}), or one not yet packed (\texttt{not\_ready}), is refused.
  \item \textbf{Collected.} The desk scans the bag before handing it over. A bag from another order
    (\texttt{wrong\_bag}) or a rider other than the assigned one (\texttt{wrong\_rider}) is refused and the order
    stays packed; the refusal names the crate slot holding the right bag.
\end{enumerate}

The time from packed to handed over is a fulfilment KPI (p90), so slow pick-ups show up next to slow packing.

\section{Order to dispatch (HLD 6.1)}

A paid order reaches a rider in seven steps, and stock leaves the books only at packing. That one rule keeps every
failure path cheap: before packing, any way out only gives a hold back.

\begin{enumerate}
  \item \textbf{Hold.} At checkout, Inventory holds the whole cart for 10 minutes, or none of it.
  \item \textbf{Pay.} Payment confirms the hold and publishes \texttt{order.paid}. The orchestrator creates the job
    and its pick tasks, oldest batch first, in walking order.
  \item \textbf{Claim.} A picker claims one order, or three at peak, with a conditional \texttt{UPDATE}.
  \item \textbf{Pick.} The picker walks the route and scans each unit into its tote.
  \item \textbf{Pack.} At the desk the sale is committed and a rider is assigned.
  \item \textbf{Check in.} The rider is told which crate slot holds the bag.
  \item \textbf{Hand over.} The desk scans the bag; only this order's bag goes to this order's rider.
\end{enumerate}

\paragraph{When an item is missing} (the HLD 9.2 saga: next bin, else substitute, else refund). Units the picker
found stay picked. The rest move to another bin with spare stock if there is one. Otherwise they are dropped and
refunded through \texttt{order.line\_dropped}, and a substitute is offered.

To accept a substitute, the demo has the picker approve it on the customer's behalf; a real app would ask the
customer on their phone. Only the item on offer can be accepted, and it is never dearer than the missing one, so
whoever accepts can never make the customer pay more than they were refunded. Only units the order really dropped
can be replaced, and they are marked so they cannot be
replaced twice. The substitute is held in the same transaction that changes the order line
(\texttt{inventory.swap\_line}, from Member~1); if it sold out since the offer, nothing changes. It joins the walk as
a new line, and \texttt{order.line\_substituted} charges for it, since the dropped units were already refunded.

\paragraph{Cancel and refunds.} Until packing, a cancel gives back every hold, and \texttt{order.cancelled} refunds
only those units. Units dropped earlier were refunded by their own \texttt{order.line\_dropped}, so the two never add up
to more than the order: in the tests, Rs~400 dropped plus Rs~560 cancelled is the order's Rs~960. After packing, a
cancel would be a return, which is out of scope, so it is refused (\texttt{too\_late}).

\section{Measurements}

Batch picking halves the walking at peak. The timing run's times are compressed, so they compare batching with
one order per walk; they do not test the HLD 12 time targets.

\paragraph{Walking distance} (\texttt{scripts/compare\_picking.py}, task 2.6). The same 30 random paid orders were
picked twice. One order per walk took 30 walks and 622 steps; batches of up to 3 took 12 walks and 322 steps, 48\%
less (46--49\% on other random seeds). A walk's distance is twice the furthest bin it visits, counted in bin steps.

\paragraph{Timing run} (\texttt{scripts/simulate\_orders.py}, task 2.8). 100 random paid orders, all waiting at the
start, were picked and packed through the real API, with one walk step taking 0.05\,s (seed 7). The stock audit was
clean after every run (Table~\ref{tab:timing-run}).

\begin{table}[ht]
  \centering
  \small
  \begin{tabular}{@{}lrrrl@{}}
    \toprule
    KPI & \shortstack[r]{One order per walk\\1 picker} & \shortstack[r]{Batches\\1 picker}
        & \shortstack[r]{Batches\\3 pickers} & HLD 12 target \\
    \midrule
    Walks                       & 100      & 34       & 34       &               \\
    Time to clear the queue (s) & 113      & 55       & 19       &               \\
    Order to packed, p90 (s)    & 103.3    & 49.6     & 17.0     & $\le 120$     \\
    Order to packed, p50 (s)    & 58.1     & 28.6     & 9.9      &               \\
    Claim to packed, p90 (s)    & 1.8      & 1.9      & 1.9      &               \\
    Fill rate                   & 100\%    & 100\%    & 100\%    & $\ge 99\%$    \\
    Items per picker-hour       & 8{,}905  & 18{,}149 & 18{,}278 & $\ge 100$     \\
    \bottomrule
  \end{tabular}
  \caption{Timing run: 100 paid orders, 0.05\,s per walk step, seed 7. The time-based targets are shown for
    reference only: with walking compressed, those rows are not a pass.}
  \label{tab:timing-run}
\end{table}

\paragraph{What the simulation can and cannot tell us.} Its comparisons hold, because every run shares the same orders,
store layout and walking speed: batching cut walks from 100 to 34 and halved the time to clear the queue. It also
shows the code stays correct under load, since every order was packed and the stock audit stayed clean.

Its absolute numbers do not hold. Walking is compressed, and scanning, packing and handover take no time, so 18,000
items per picker-hour describes the code, not people. All orders arrive at once, so order-to-packed mostly measures how
fast the pickers drain the queue.

Fill rate is 100\% only because no shelf is ever empty in the simulation. Real numbers need timed pickers in a real or
mock-up store, with orders arriving through the day.

\section{Trade-off: batch picking at peak (HLD 13)}

We batch orders only when at least 4 are waiting. Batching halves the walking at peak, and the threshold keeps it from
making an order wait when the store is quiet.

HLD 13.1 chose ``single-order off-peak, batch picking at peak'' because ``batching raises throughput at peak but adds
wait when orders are few''. We measured the first half: 48\% less walking for 30 orders. With one picker, a 100-order
queue cleared in 55\,s instead of 113\,s (compressed time).

The second half is handled by design rather than measured. With fewer than 4 orders waiting, a picker takes one at
once, so no order waits for a batch to fill. With 4 or more, a batch of 3 still leaves at least one order for the next
picker.

What batching still costs:
\begin{itemize}
  \item Each order waits while the picker walks to the other orders' bins. In the timing run, claim-to-packed p90
    barely moved (1.8\,s to 1.9\,s), but on a real floor this is real walking time.
  \item The picker must put each item in the right tote, and nothing checks the tote yet.
  \item A batch is the 3 oldest orders, not 3 whose items sit close together, so some batches save little walking.
\end{itemize}

\section{Limits and open items}

The fulfilment prototype's limits are shortcuts, each with a clear next step (Table~\ref{tab:limits}).

\begin{table}[!htb]
  \centering
  \footnotesize
  \begin{tabular}{@{}>{\raggedright\arraybackslash}p{0.27\linewidth}>{\raggedright\arraybackslash}p{0.36\linewidth}
                  >{\raggedright\arraybackslash}p{0.31\linewidth}@{}}
    \toprule
    Limit & What it means & Next step \\
    \midrule
    The picker accepts a substitute for the customer
      & The customer is not asked; the picker decides on their behalf, but only for the offered item, never a dearer one
      & Ask the customer on their phone, with a timeout \\
    \addlinespace[3pt]
    Units already in a tote when their batch expires
      & They are still packed and delivered; only the rest of the batch is refused at the scanner
      & Check the expiry date again at packing \\
    \addlinespace[3pt]
    Units typed as ``found'' are not scanned
      & They skip the barcode check
      & Scan them before packing \\
    \addlinespace[3pt]
    The picker check is opt-in
      & A handheld that sends no \texttt{picker\_id} can still act on an order since given to someone else
      & The demo station and phone screen send it; make it required \\
    \addlinespace[3pt]
    Riders are made up
      & A fixed list stands in for the rider system, and the desk checks riders in
      & Connect a real rider system \\
    \addlinespace[3pt]
    Refunds use today's price
      & A price change between checkout and refund would refund the wrong amount
      & Store the price on each reservation line \\
    \addlinespace[3pt]
    Batches are the oldest orders
      & Orders are not grouped by where their items sit
      & Group by zone at peak \\
    \bottomrule
  \end{tabular}
  \caption{Limits of the fulfilment prototype and what would close them.}
  \label{tab:limits}
\end{table}
. Needs: array, booktabs, tikz (arrows.meta,
% positioning, calc).
% member2-standalone.tex compiles these sections on their own.

\section{Fulfilment Orchestrator (HLD 5, service 4)}

The Fulfilment Orchestrator turns each paid order into a job and keeps it moving until a rider takes the bag, so that
no order gets stuck or is acted on twice. It listens for \texttt{order.paid}, creates a \texttt{FulfilmentJob} with
one pick task per bin (oldest batch first, in walking order) and owns the job's state.

Every state change is one conditional \texttt{UPDATE} that names the state it expects:
\begin{center}
  \texttt{UPDATE fulfilment\_jobs SET state = :new WHERE id = :job AND state = :expected}
\end{center}
Of two requests racing for the same change, exactly one wins and the other is told what happened. This covers
claiming, cancelling, scanning, packing, stall release and handover.

\begin{figure}[ht]
  \centering
  \begin{tikzpicture}[
      >={Stealth[length=2.2mm]}, node distance=17mm,
      state/.style={draw=black!45, rounded corners=3pt, minimum width=24mm, minimum height=10mm, font=\bfseries},
      sale/.style={state, draw=blue!70!black, fill=blue!8, line width=0.9pt},
      edgelabel/.style={font=\small, text=black!65}]
    \node[state] (created) {created};
    \node[state, right=of created] (picking) {picking};
    \node[sale, right=of picking] (packed) {packed};
    \node[state, right=of packed] (handed) {handed over};
    \node[state, below=16mm of $(created)!0.5!(picking)$] (cancelled) {cancelled};
    \node[edgelabel, below=1mm of packed] {sale committed here};

    \draw[->, black!55] (created) -- node[edgelabel, above] {claim} (picking);
    \draw[->, black!55] (picking) -- node[edgelabel, above] {pack} (packed);
    \draw[->, black!55] (packed) -- node[edgelabel, above] {scan bag} (handed);
    \draw[->, black!55] (picking.north) -- ++(0, 7mm) -| node[edgelabel, pos=0.25, above]
      {idle 5 min, nothing in tote} (created.north);
    \draw[->, black!55] (created.south) |- node[edgelabel, pos=0.3, left] {cancel} (cancelled.west);
    \draw[->, black!55] (picking.south) |- node[edgelabel, pos=0.3, right] {cancel, or nothing to pack}
      (cancelled.east);
  \end{tikzpicture}
  \caption{Fulfilment job lifecycle. Stock leaves the books only at packing; every earlier exit only releases holds.}
  \label{fig:job-lifecycle}
\end{figure}

Figure~\ref{fig:job-lifecycle} shows the lifecycle. Before packing, an order can leave the line by a cancel, or go back
to the queue when its picker goes idle; both only release holds. After packing, the only way on is to its own rider.

\begin{itemize}
  \item \textbf{Store load.} More than 5 orders waiting marks the store \texttt{busy}, and serviceability adds
    5 minutes to the ETA.
  \item \textbf{Cancel after payment.} Allowed while the order is \texttt{created} or \texttt{picking}: its holds are
    given back and nothing leaves on-hand stock. A packed order would be a return, which is out of scope.
  \item \textbf{Stuck orders.} No claim, scan or shelf report for 5 minutes: with nothing in the tote, the order goes
    back to the queue (checked every 60\,s); with items in a tote, it is flagged for a person. Packed bags waiting
    over 5 minutes for a rider are flagged too.
  \item \textbf{KPIs} over the last 30 days: order-to-packed p50 and p90, claim-to-packed p90, packed-to-handover
    p90, fill rate (units picked $\div$ units ordered on the reservation) and items per picker-hour (a picker's
    overlapping orders counted once).
\end{itemize}

Percentiles take the value at index $\lfloor q \times (n - 1) \rfloor$ of the sorted list, the nearest rank below.
With few orders this errs low: the p90 of 30, 60 and 90\,s is 60\,s.

\section{Pick \& Pack (HLD 5, service 5)}

Pick \& Pack gives pickers work, checks every scan and packs the order. Packing is the moment the sale leaves the
books; until then the units are only held, so a cancel never has to put stock back.

\begin{itemize}
  \item \textbf{Claiming.} Off-peak, a picker gets the oldest waiting order. With 4 or more waiting, they get the
    3 oldest, each in its own tote (A, B, C), with one route through all their lines in walking order.
  \item \textbf{Scanning.} Every unit is barcode-checked as it goes into the tote (``3 of 6 scanned''). A wrong item,
    or one unit more than the line needs, changes nothing.
  \item \textbf{Empty shelf.} The picker says how many were there, and those stay picked. For the rest, the count is
    corrected (a \texttt{not\_found} ledger row) and the line moves to another bin with spare stock, never units
    another customer holds. With none left it is dropped: the hold is released, \texttt{order.line\_dropped} carries
    the refund, and a substitute is offered (same category, enough to sell, no dearer than the missing item, nearest
    price). The rest of the bin is written off too, so customers cannot buy stock that is not there. Units another
    unpacked order is waiting for stay, so held units never exceed those on hand, until that order's picker reports
    the same empty bin.
  \item \textbf{Packing.} Commits the sale for every picked line (\texttt{qty\_on\_hand} and \texttt{qty\_reserved}
    both drop) and publishes \texttt{order.ready\_for\_pickup}. An order with nothing picked is cancelled instead.
  \item \textbf{The handheld.} \texttt{/picker.html} is a phone screen: one line at a time (``Go to bin A-01-04,
    pick 2 $\times$ Basmati rice''), big buttons, and the tote shown at peak. A handheld that sends its picker id
    cannot act on an order that has since gone to someone else.
\end{itemize}

\section{Rider Handoff (HLD 5, service 11)}

The Rider Handoff Adapter makes sure the right bag goes to the right rider. Before it, handover was one unchecked
call, and a rider could leave with someone else's order.

\begin{enumerate}
  \item \textbf{Assigned.} When an order is packed, the rider system sends a rider. In the prototype a stand-in picks
    one from a fixed list (R-11 to R-15) inside the packing transaction, so a packed order always has a rider.
  \item \textbf{Arrived.} The rider checks in at the desk for an order and is told its crate slot. Another rider's
    order (\texttt{wrong\_rider}), or one not yet packed (\texttt{not\_ready}), is refused.
  \item \textbf{Collected.} The desk scans the bag before handing it over. A bag from another order
    (\texttt{wrong\_bag}) or a rider other than the assigned one (\texttt{wrong\_rider}) is refused and the order
    stays packed; the refusal names the crate slot holding the right bag.
\end{enumerate}

The time from packed to handed over is a fulfilment KPI (p90), so slow pick-ups show up next to slow packing.

\section{Order to dispatch (HLD 6.1)}

A paid order reaches a rider in seven steps, and stock leaves the books only at packing. That one rule keeps every
failure path cheap: before packing, any way out only gives a hold back.

\begin{enumerate}
  \item \textbf{Hold.} At checkout, Inventory holds the whole cart for 10 minutes, or none of it.
  \item \textbf{Pay.} Payment confirms the hold and publishes \texttt{order.paid}. The orchestrator creates the job
    and its pick tasks, oldest batch first, in walking order.
  \item \textbf{Claim.} A picker claims one order, or three at peak, with a conditional \texttt{UPDATE}.
  \item \textbf{Pick.} The picker walks the route and scans each unit into its tote.
  \item \textbf{Pack.} At the desk the sale is committed and a rider is assigned.
  \item \textbf{Check in.} The rider is told which crate slot holds the bag.
  \item \textbf{Hand over.} The desk scans the bag; only this order's bag goes to this order's rider.
\end{enumerate}

\paragraph{When an item is missing} (the HLD 9.2 saga: next bin, else substitute, else refund). Units the picker
found stay picked. The rest move to another bin with spare stock if there is one. Otherwise they are dropped and
refunded through \texttt{order.line\_dropped}, and a substitute is offered.

To accept a substitute, the demo has the picker approve it on the customer's behalf; a real app would ask the
customer on their phone. Only the item on offer can be accepted, and it is never dearer than the missing one, so
whoever accepts can never make the customer pay more than they were refunded. Only units the order really dropped
can be replaced, and they are marked so they cannot be
replaced twice. The substitute is held in the same transaction that changes the order line
(\texttt{inventory.swap\_line}, from Member~1); if it sold out since the offer, nothing changes. It joins the walk as
a new line, and \texttt{order.line\_substituted} charges for it, since the dropped units were already refunded.

\paragraph{Cancel and refunds.} Until packing, a cancel gives back every hold, and \texttt{order.cancelled} refunds
only those units. Units dropped earlier were refunded by their own \texttt{order.line\_dropped}, so the two never add up
to more than the order: in the tests, Rs~400 dropped plus Rs~560 cancelled is the order's Rs~960. After packing, a
cancel would be a return, which is out of scope, so it is refused (\texttt{too\_late}).

\section{Measurements}

Batch picking halves the walking at peak. The timing run's times are compressed, so they compare batching with
one order per walk; they do not test the HLD 12 time targets.

\paragraph{Walking distance} (\texttt{scripts/compare\_picking.py}, task 2.6). The same 30 random paid orders were
picked twice. One order per walk took 30 walks and 622 steps; batches of up to 3 took 12 walks and 322 steps, 48\%
less (46--49\% on other random seeds). A walk's distance is twice the furthest bin it visits, counted in bin steps.

\paragraph{Timing run} (\texttt{scripts/simulate\_orders.py}, task 2.8). 100 random paid orders, all waiting at the
start, were picked and packed through the real API, with one walk step taking 0.05\,s (seed 7). The stock audit was
clean after every run (Table~\ref{tab:timing-run}).

\begin{table}[ht]
  \centering
  \small
  \begin{tabular}{@{}lrrrl@{}}
    \toprule
    KPI & \shortstack[r]{One order per walk\\1 picker} & \shortstack[r]{Batches\\1 picker}
        & \shortstack[r]{Batches\\3 pickers} & HLD 12 target \\
    \midrule
    Walks                       & 100      & 34       & 34       &               \\
    Time to clear the queue (s) & 113      & 55       & 19       &               \\
    Order to packed, p90 (s)    & 103.3    & 49.6     & 17.0     & $\le 120$     \\
    Order to packed, p50 (s)    & 58.1     & 28.6     & 9.9      &               \\
    Claim to packed, p90 (s)    & 1.8      & 1.9      & 1.9      &               \\
    Fill rate                   & 100\%    & 100\%    & 100\%    & $\ge 99\%$    \\
    Items per picker-hour       & 8{,}905  & 18{,}149 & 18{,}278 & $\ge 100$     \\
    \bottomrule
  \end{tabular}
  \caption{Timing run: 100 paid orders, 0.05\,s per walk step, seed 7. The time-based targets are shown for
    reference only: with walking compressed, those rows are not a pass.}
  \label{tab:timing-run}
\end{table}

\paragraph{What the simulation can and cannot tell us.} Its comparisons hold, because every run shares the same orders,
store layout and walking speed: batching cut walks from 100 to 34 and halved the time to clear the queue. It also
shows the code stays correct under load, since every order was packed and the stock audit stayed clean.

Its absolute numbers do not hold. Walking is compressed, and scanning, packing and handover take no time, so 18,000
items per picker-hour describes the code, not people. All orders arrive at once, so order-to-packed mostly measures how
fast the pickers drain the queue.

Fill rate is 100\% only because no shelf is ever empty in the simulation. Real numbers need timed pickers in a real or
mock-up store, with orders arriving through the day.

\section{Trade-off: batch picking at peak (HLD 13)}

We batch orders only when at least 4 are waiting. Batching halves the walking at peak, and the threshold keeps it from
making an order wait when the store is quiet.

HLD 13.1 chose ``single-order off-peak, batch picking at peak'' because ``batching raises throughput at peak but adds
wait when orders are few''. We measured the first half: 48\% less walking for 30 orders. With one picker, a 100-order
queue cleared in 55\,s instead of 113\,s (compressed time).

The second half is handled by design rather than measured. With fewer than 4 orders waiting, a picker takes one at
once, so no order waits for a batch to fill. With 4 or more, a batch of 3 still leaves at least one order for the next
picker.

What batching still costs:
\begin{itemize}
  \item Each order waits while the picker walks to the other orders' bins. In the timing run, claim-to-packed p90
    barely moved (1.8\,s to 1.9\,s), but on a real floor this is real walking time.
  \item The picker must put each item in the right tote, and nothing checks the tote yet.
  \item A batch is the 3 oldest orders, not 3 whose items sit close together, so some batches save little walking.
\end{itemize}

\section{Limits and open items}

The fulfilment prototype's limits are shortcuts, each with a clear next step (Table~\ref{tab:limits}).

\begin{table}[!htb]
  \centering
  \footnotesize
  \begin{tabular}{@{}>{\raggedright\arraybackslash}p{0.27\linewidth}>{\raggedright\arraybackslash}p{0.36\linewidth}
                  >{\raggedright\arraybackslash}p{0.31\linewidth}@{}}
    \toprule
    Limit & What it means & Next step \\
    \midrule
    The picker accepts a substitute for the customer
      & The customer is not asked; the picker decides on their behalf, but only for the offered item, never a dearer one
      & Ask the customer on their phone, with a timeout \\
    \addlinespace[3pt]
    Units already in a tote when their batch expires
      & They are still packed and delivered; only the rest of the batch is refused at the scanner
      & Check the expiry date again at packing \\
    \addlinespace[3pt]
    Units typed as ``found'' are not scanned
      & They skip the barcode check
      & Scan them before packing \\
    \addlinespace[3pt]
    The picker check is opt-in
      & A handheld that sends no \texttt{picker\_id} can still act on an order since given to someone else
      & The demo station and phone screen send it; make it required \\
    \addlinespace[3pt]
    Riders are made up
      & A fixed list stands in for the rider system, and the desk checks riders in
      & Connect a real rider system \\
    \addlinespace[3pt]
    Refunds use today's price
      & A price change between checkout and refund would refund the wrong amount
      & Store the price on each reservation line \\
    \addlinespace[3pt]
    Batches are the oldest orders
      & Orders are not grouped by where their items sit
      & Group by zone at peak \\
    \bottomrule
  \end{tabular}
  \caption{Limits of the fulfilment prototype and what would close them.}
  \label{tab:limits}
\end{table}
. Needs: array, booktabs, tikz (arrows.meta,
% positioning, calc).
% member2-standalone.tex compiles these sections on their own.

\section{Fulfilment Orchestrator (HLD 5, service 4)}

The Fulfilment Orchestrator turns each paid order into a job and keeps it moving until a rider takes the bag, so that
no order gets stuck or is acted on twice. It listens for \texttt{order.paid}, creates a \texttt{FulfilmentJob} with
one pick task per bin (oldest batch first, in walking order) and owns the job's state.

Every state change is one conditional \texttt{UPDATE} that names the state it expects:
\begin{center}
  \texttt{UPDATE fulfilment\_jobs SET state = :new WHERE id = :job AND state = :expected}
\end{center}
Of two requests racing for the same change, exactly one wins and the other is told what happened. This covers
claiming, cancelling, scanning, packing, stall release and handover.

\begin{figure}[ht]
  \centering
  \begin{tikzpicture}[
      >={Stealth[length=2.2mm]}, node distance=17mm,
      state/.style={draw=black!45, rounded corners=3pt, minimum width=24mm, minimum height=10mm, font=\bfseries},
      sale/.style={state, draw=blue!70!black, fill=blue!8, line width=0.9pt},
      edgelabel/.style={font=\small, text=black!65}]
    \node[state] (created) {created};
    \node[state, right=of created] (picking) {picking};
    \node[sale, right=of picking] (packed) {packed};
    \node[state, right=of packed] (handed) {handed over};
    \node[state, below=16mm of $(created)!0.5!(picking)$] (cancelled) {cancelled};
    \node[edgelabel, below=1mm of packed] {sale committed here};

    \draw[->, black!55] (created) -- node[edgelabel, above] {claim} (picking);
    \draw[->, black!55] (picking) -- node[edgelabel, above] {pack} (packed);
    \draw[->, black!55] (packed) -- node[edgelabel, above] {scan bag} (handed);
    \draw[->, black!55] (picking.north) -- ++(0, 7mm) -| node[edgelabel, pos=0.25, above]
      {idle 5 min, nothing in tote} (created.north);
    \draw[->, black!55] (created.south) |- node[edgelabel, pos=0.3, left] {cancel} (cancelled.west);
    \draw[->, black!55] (picking.south) |- node[edgelabel, pos=0.3, right] {cancel, or nothing to pack}
      (cancelled.east);
  \end{tikzpicture}
  \caption{Fulfilment job lifecycle. Stock leaves the books only at packing; every earlier exit only releases holds.}
  \label{fig:job-lifecycle}
\end{figure}

Figure~\ref{fig:job-lifecycle} shows the lifecycle. Before packing, an order can leave the line by a cancel, or go back
to the queue when its picker goes idle; both only release holds. After packing, the only way on is to its own rider.

\begin{itemize}
  \item \textbf{Store load.} More than 5 orders waiting marks the store \texttt{busy}, and serviceability adds
    5 minutes to the ETA.
  \item \textbf{Cancel after payment.} Allowed while the order is \texttt{created} or \texttt{picking}: its holds are
    given back and nothing leaves on-hand stock. A packed order would be a return, which is out of scope.
  \item \textbf{Stuck orders.} No claim, scan or shelf report for 5 minutes: with nothing in the tote, the order goes
    back to the queue (checked every 60\,s); with items in a tote, it is flagged for a person. Packed bags waiting
    over 5 minutes for a rider are flagged too.
  \item \textbf{KPIs} over the last 30 days: order-to-packed p50 and p90, claim-to-packed p90, packed-to-handover
    p90, fill rate (units picked $\div$ units ordered on the reservation) and items per picker-hour (a picker's
    overlapping orders counted once).
\end{itemize}

Percentiles take the value at index $\lfloor q \times (n - 1) \rfloor$ of the sorted list, the nearest rank below.
With few orders this errs low: the p90 of 30, 60 and 90\,s is 60\,s.

\section{Pick \& Pack (HLD 5, service 5)}

Pick \& Pack gives pickers work, checks every scan and packs the order. Packing is the moment the sale leaves the
books; until then the units are only held, so a cancel never has to put stock back.

\begin{itemize}
  \item \textbf{Claiming.} Off-peak, a picker gets the oldest waiting order. With 4 or more waiting, they get the
    3 oldest, each in its own tote (A, B, C), with one route through all their lines in walking order.
  \item \textbf{Scanning.} Every unit is barcode-checked as it goes into the tote (``3 of 6 scanned''). A wrong item,
    or one unit more than the line needs, changes nothing.
  \item \textbf{Empty shelf.} The picker says how many were there, and those stay picked. For the rest, the count is
    corrected (a \texttt{not\_found} ledger row) and the line moves to another bin with spare stock, never units
    another customer holds. With none left it is dropped: the hold is released, \texttt{order.line\_dropped} carries
    the refund, and a substitute is offered (same category, enough to sell, no dearer than the missing item, nearest
    price). The rest of the bin is written off too, so customers cannot buy stock that is not there. Units another
    unpacked order is waiting for stay, so held units never exceed those on hand, until that order's picker reports
    the same empty bin.
  \item \textbf{Packing.} Commits the sale for every picked line (\texttt{qty\_on\_hand} and \texttt{qty\_reserved}
    both drop) and publishes \texttt{order.ready\_for\_pickup}. An order with nothing picked is cancelled instead.
  \item \textbf{The handheld.} \texttt{/picker.html} is a phone screen: one line at a time (``Go to bin A-01-04,
    pick 2 $\times$ Basmati rice''), big buttons, and the tote shown at peak. A handheld that sends its picker id
    cannot act on an order that has since gone to someone else.
\end{itemize}

\section{Rider Handoff (HLD 5, service 11)}

The Rider Handoff Adapter makes sure the right bag goes to the right rider. Before it, handover was one unchecked
call, and a rider could leave with someone else's order.

\begin{enumerate}
  \item \textbf{Assigned.} When an order is packed, the rider system sends a rider. In the prototype a stand-in picks
    one from a fixed list (R-11 to R-15) inside the packing transaction, so a packed order always has a rider.
  \item \textbf{Arrived.} The rider checks in at the desk for an order and is told its crate slot. Another rider's
    order (\texttt{wrong\_rider}), or one not yet packed (\texttt{not\_ready}), is refused.
  \item \textbf{Collected.} The desk scans the bag before handing it over. A bag from another order
    (\texttt{wrong\_bag}) or a rider other than the assigned one (\texttt{wrong\_rider}) is refused and the order
    stays packed; the refusal names the crate slot holding the right bag.
\end{enumerate}

The time from packed to handed over is a fulfilment KPI (p90), so slow pick-ups show up next to slow packing.

\section{Order to dispatch (HLD 6.1)}

A paid order reaches a rider in seven steps, and stock leaves the books only at packing. That one rule keeps every
failure path cheap: before packing, any way out only gives a hold back.

\begin{enumerate}
  \item \textbf{Hold.} At checkout, Inventory holds the whole cart for 10 minutes, or none of it.
  \item \textbf{Pay.} Payment confirms the hold and publishes \texttt{order.paid}. The orchestrator creates the job
    and its pick tasks, oldest batch first, in walking order.
  \item \textbf{Claim.} A picker claims one order, or three at peak, with a conditional \texttt{UPDATE}.
  \item \textbf{Pick.} The picker walks the route and scans each unit into its tote.
  \item \textbf{Pack.} At the desk the sale is committed and a rider is assigned.
  \item \textbf{Check in.} The rider is told which crate slot holds the bag.
  \item \textbf{Hand over.} The desk scans the bag; only this order's bag goes to this order's rider.
\end{enumerate}

\paragraph{When an item is missing} (the HLD 9.2 saga: next bin, else substitute, else refund). Units the picker
found stay picked. The rest move to another bin with spare stock if there is one. Otherwise they are dropped and
refunded through \texttt{order.line\_dropped}, and a substitute is offered.

To accept a substitute, the demo has the picker approve it on the customer's behalf; a real app would ask the
customer on their phone. Only the item on offer can be accepted, and it is never dearer than the missing one, so
whoever accepts can never make the customer pay more than they were refunded. Only units the order really dropped
can be replaced, and they are marked so they cannot be
replaced twice. The substitute is held in the same transaction that changes the order line
(\texttt{inventory.swap\_line}, from Member~1); if it sold out since the offer, nothing changes. It joins the walk as
a new line, and \texttt{order.line\_substituted} charges for it, since the dropped units were already refunded.

\paragraph{Cancel and refunds.} Until packing, a cancel gives back every hold, and \texttt{order.cancelled} refunds
only those units. Units dropped earlier were refunded by their own \texttt{order.line\_dropped}, so the two never add up
to more than the order: in the tests, Rs~400 dropped plus Rs~560 cancelled is the order's Rs~960. After packing, a
cancel would be a return, which is out of scope, so it is refused (\texttt{too\_late}).

\section{Measurements}

Batch picking halves the walking at peak. The timing run's times are compressed, so they compare batching with
one order per walk; they do not test the HLD 12 time targets.

\paragraph{Walking distance} (\texttt{scripts/compare\_picking.py}, task 2.6). The same 30 random paid orders were
picked twice. One order per walk took 30 walks and 622 steps; batches of up to 3 took 12 walks and 322 steps, 48\%
less (46--49\% on other random seeds). A walk's distance is twice the furthest bin it visits, counted in bin steps.

\paragraph{Timing run} (\texttt{scripts/simulate\_orders.py}, task 2.8). 100 random paid orders, all waiting at the
start, were picked and packed through the real API, with one walk step taking 0.05\,s (seed 7). The stock audit was
clean after every run (Table~\ref{tab:timing-run}).

\begin{table}[ht]
  \centering
  \small
  \begin{tabular}{@{}lrrrl@{}}
    \toprule
    KPI & \shortstack[r]{One order per walk\\1 picker} & \shortstack[r]{Batches\\1 picker}
        & \shortstack[r]{Batches\\3 pickers} & HLD 12 target \\
    \midrule
    Walks                       & 100      & 34       & 34       &               \\
    Time to clear the queue (s) & 113      & 55       & 19       &               \\
    Order to packed, p90 (s)    & 103.3    & 49.6     & 17.0     & $\le 120$     \\
    Order to packed, p50 (s)    & 58.1     & 28.6     & 9.9      &               \\
    Claim to packed, p90 (s)    & 1.8      & 1.9      & 1.9      &               \\
    Fill rate                   & 100\%    & 100\%    & 100\%    & $\ge 99\%$    \\
    Items per picker-hour       & 8{,}905  & 18{,}149 & 18{,}278 & $\ge 100$     \\
    \bottomrule
  \end{tabular}
  \caption{Timing run: 100 paid orders, 0.05\,s per walk step, seed 7. The time-based targets are shown for
    reference only: with walking compressed, those rows are not a pass.}
  \label{tab:timing-run}
\end{table}

\paragraph{What the simulation can and cannot tell us.} Its comparisons hold, because every run shares the same orders,
store layout and walking speed: batching cut walks from 100 to 34 and halved the time to clear the queue. It also
shows the code stays correct under load, since every order was packed and the stock audit stayed clean.

Its absolute numbers do not hold. Walking is compressed, and scanning, packing and handover take no time, so 18,000
items per picker-hour describes the code, not people. All orders arrive at once, so order-to-packed mostly measures how
fast the pickers drain the queue.

Fill rate is 100\% only because no shelf is ever empty in the simulation. Real numbers need timed pickers in a real or
mock-up store, with orders arriving through the day.

\section{Trade-off: batch picking at peak (HLD 13)}

We batch orders only when at least 4 are waiting. Batching halves the walking at peak, and the threshold keeps it from
making an order wait when the store is quiet.

HLD 13.1 chose ``single-order off-peak, batch picking at peak'' because ``batching raises throughput at peak but adds
wait when orders are few''. We measured the first half: 48\% less walking for 30 orders. With one picker, a 100-order
queue cleared in 55\,s instead of 113\,s (compressed time).

The second half is handled by design rather than measured. With fewer than 4 orders waiting, a picker takes one at
once, so no order waits for a batch to fill. With 4 or more, a batch of 3 still leaves at least one order for the next
picker.

What batching still costs:
\begin{itemize}
  \item Each order waits while the picker walks to the other orders' bins. In the timing run, claim-to-packed p90
    barely moved (1.8\,s to 1.9\,s), but on a real floor this is real walking time.
  \item The picker must put each item in the right tote, and nothing checks the tote yet.
  \item A batch is the 3 oldest orders, not 3 whose items sit close together, so some batches save little walking.
\end{itemize}

\section{Limits and open items}

The fulfilment prototype's limits are shortcuts, each with a clear next step (Table~\ref{tab:limits}).

\begin{table}[!htb]
  \centering
  \footnotesize
  \begin{tabular}{@{}>{\raggedright\arraybackslash}p{0.27\linewidth}>{\raggedright\arraybackslash}p{0.36\linewidth}
                  >{\raggedright\arraybackslash}p{0.31\linewidth}@{}}
    \toprule
    Limit & What it means & Next step \\
    \midrule
    The picker accepts a substitute for the customer
      & The customer is not asked; the picker decides on their behalf, but only for the offered item, never a dearer one
      & Ask the customer on their phone, with a timeout \\
    \addlinespace[3pt]
    Units already in a tote when their batch expires
      & They are still packed and delivered; only the rest of the batch is refused at the scanner
      & Check the expiry date again at packing \\
    \addlinespace[3pt]
    Units typed as ``found'' are not scanned
      & They skip the barcode check
      & Scan them before packing \\
    \addlinespace[3pt]
    The picker check is opt-in
      & A handheld that sends no \texttt{picker\_id} can still act on an order since given to someone else
      & The demo station and phone screen send it; make it required \\
    \addlinespace[3pt]
    Riders are made up
      & A fixed list stands in for the rider system, and the desk checks riders in
      & Connect a real rider system \\
    \addlinespace[3pt]
    Refunds use today's price
      & A price change between checkout and refund would refund the wrong amount
      & Store the price on each reservation line \\
    \addlinespace[3pt]
    Batches are the oldest orders
      & Orders are not grouped by where their items sit
      & Group by zone at peak \\
    \bottomrule
  \end{tabular}
  \caption{Limits of the fulfilment prototype and what would close them.}
  \label{tab:limits}
\end{table}
